\documentclass[11pt,a4paper]{article}

\usepackage[utf8]{inputenc}
\usepackage[T2A]{fontenc}
\usepackage[russian]{babel}
\usepackage[left=1.7cm,right=1.7cm,top=2.0cm,bottom=1.7cm]{geometry}
\usepackage{amsmath,amssymb}
\usepackage{xcolor}
\usepackage{tikz}
\usepackage[most]{tcolorbox}
\usepackage{booktabs}
\usepackage{hyperref}
\usepackage{enumitem}
\usepackage{fancyhdr}
\usepackage{titlesec}
\usepackage{needspace}
\usepackage{microtype}

\usetikzlibrary{arrows.meta,positioning}

\definecolor{ink}{RGB}{28,28,30}
\definecolor{paper}{RGB}{252,250,246}
\definecolor{rulec}{RGB}{200,190,178}
\definecolor{mvred}{RGB}{196,30,38}
\definecolor{inksoft}{RGB}{95,92,88}
\definecolor{mint}{RGB}{40,105,95}
\definecolor{sand}{RGB}{236,226,210}
\definecolor{slate}{RGB}{55,70,85}

\hypersetup{colorlinks=true,linkcolor=mvred,urlcolor=mvred}

\pagestyle{fancy}
\fancyhf{}
\renewcommand{\headrulewidth}{0.25pt}
\fancyhead[L]{\small\textcolor{inksoft}{М.Видео $\cdot$ Intelligent Search}}
\fancyhead[R]{\small\textcolor{inksoft}{Roadmap}}
\fancyfoot[C]{\thepage}
\setlength{\headheight}{14pt}

\titleformat{\section}{\large\bfseries\color{ink}}{\thesection.}{0.4em}{}
\titleformat{\subsection}{\normalsize\bfseries\color{slate}}{\thesubsection.}{0.35em}{}
\titlespacing*{\section}{0pt}{1.15em}{0.65em}

\newcommand{\secrule}{%
  \par\vspace{4pt}%
  {\color{mvred}\noindent\rule{\linewidth}{0.9pt}\par}%
  \vspace{8pt}%
}

% Keep section title + intro + card together (no orphan headers / blank tails)
\newcommand{\daysection}[1]{%
  \Needspace{22\baselineskip}%
  \section{#1}\secrule
}

\setlist[enumerate]{leftmargin=1.3em,itemsep=0.18em,topsep=0.2em,parsep=0pt}
\setlength{\parskip}{0.3em}
\setlength{\parindent}{0pt}

% Non-breakable card: either fully on page, or whole block moves down via Needspace
\newcommand{\dayblock}[4]{%
  \begin{tcolorbox}[
    enhanced,
    colback=white,
    colframe=rulec,
    boxrule=0.65pt,
    arc=2pt,
    left=8pt,
    right=8pt,
    top=8pt,
    bottom=8pt,
    before skip=3pt,
    after skip=11pt,
  ]
  {\color{mvred}\bfseries #1}~~{\bfseries #2}\par\vspace{4pt}%
  {\scriptsize\bfseries\color{inksoft}ЧТО ДЕЛАЕМ}\par\vspace{1pt}%
  #3\par\vspace{5pt}%
  {\scriptsize\bfseries\color{mint}ЧТО ДОЛЖНО ПОЛУЧИТЬСЯ}\par\vspace{1pt}%
  #4%
  \end{tcolorbox}%
}

\begin{document}
\pagecolor{paper}
\color{ink}

\noindent{\footnotesize\bfseries\color{inksoft}БУТКЕМП $\cdot$ КЕЙС М.ВИДЕО $\cdot$ NLP / SEARCH}

\vspace{0.35cm}
{\fontsize{24}{28}\selectfont\bfseries Roadmap}\\[0.08cm]
{\fontsize{13}{16}\selectfont\color{mvred}\bfseries от теории к standalone-поисковику}

\vspace{0.28cm}
{\color{mvred}\hrule height 1.1pt}
\vspace{0.32cm}

Сначала читаем и договариваемся, о чём задача.
К середине буткемпа --- EDA и живой MVP.
Дальше своя разметка, жёстче факты, поиск по эмбеддингам и упаковка под защиту.

\vspace{0.3cm}
\noindent
\begin{tabular}{@{}p{0.47\textwidth}@{\hspace{0.06\textwidth}}p{0.47\textwidth}@{}}
{\scriptsize\bfseries\color{inksoft}ПРИНЦИП} &
{\scriptsize\bfseries\color{inksoft}СРОКИ}\\[2pt]
Факты из запроса должны помогать найти SKU, а не просто красиво лежать в JSON. &
EDA + MVP к дню 2--3. Параллельно копим свою разметку и не врём себе на метриках.
\end{tabular}

\vspace{0.4cm}
\section*{Как это читать}
\secrule
План кривоватый по нагрузке --- так и задумано.
В начале теория, в середине инженерия, ближе к защите --- продукт и история для стейкхолдеров.

\vspace{0.2cm}
\begin{center}
\begin{tikzpicture}[
  box/.style={draw=rulec, fill=white, inner sep=4pt, font=\scriptsize\bfseries, rounded corners=1.5pt},
  arr/.style={-{Stealth[length=2mm]}, thick, inksoft}
]
  \node[box] (t) {0 теория};
  \node[box, right=0.2 of t] (a) {1 каркас};
  \node[box, right=0.2 of a] (b) {2 EDA+MVP};
  \node[box, right=0.2 of b] (c) {3 разметка};
  \node[box, right=0.2 of c] (d) {4 поиск};
  \node[box, right=0.2 of d] (e) {5 защита};
  \draw[arr] (t)--(a); \draw[arr] (a)--(b); \draw[arr] (b)--(c);
  \draw[arr] (c)--(d); \draw[arr] (d)--(e);
\end{tikzpicture}\\[3pt]
{\tiny\color{mvred}под <<2 EDA+MVP>> --- жёсткий чекпоинт}
\end{center}

\daysection{День 0 --- до буткемпа: просто понять язык}
Пока никто не кодит ради скрина. Читаем и договариваемся, что значат слова.

\dayblock{00}{Теория и насмотренность}{%
\begin{enumerate}
  \item Лонгриды и семинары: признаки, overfitting, train/test, метрики.
  \item NLP: токены, BIO, чем NER отличается от классификации всего запроса.
  \item CRF vs <<сразу нейросеть>>: когда хватает модели на признаках токена.
  \item Эмбеддинги и косинус --- задел под поиск, не обязательно в день 1.
  \item Трансформеры понимать ок, но тяжёлый BERT в MVP под SLA $<100$ мс не тащим.
\end{enumerate}
}{%
\begin{enumerate}
  \item Общий язык: факт, SKU, weak label, entity F1, latency.
  \item Короткий конспект: что прочитали и что берём в кейс.
\end{enumerate}
}
\noindent\textit{Читать:} лонгриды курса, последовательности / attention, TF-IDF, $P/R/F1$.

\daysection{День 01 --- Погружение, но сразу лезем в данные}
День не заканчивается планом на доске. К вечеру parquet уже открыт.

\dayblock{01}{Синхронизация + репо + первый взгляд}{%
\begin{enumerate}
  \item Сверить кейс: запрос $\to$ JSON с фактами, быстрее 100 мс.
  \item Роли набросать (данные, модель, сервис/UI, доки).
  \item Зафиксировать JSON: brand, category, attributes, entities, latency\_ms.
  \item Репозиторий, сырые гигабайты в gitignore, черновик README.
  \item Открыть query\_clicks и sku\_desc: схемы, размеры, штук 20 строк.
\end{enumerate}
}{%
\begin{enumerate}
  \item Этот roadmap --- план на дни 2--5.
  \item Ноутбук overview: колонки, объёмы.
\end{enumerate}
}

\daysection{День 02 --- Сюда нельзя опоздать: EDA + MVP}
Если к вечеру нет JSON --- на день 3 не прыгаем в приложение мечты.

\dayblock{02}{EDA, словари, CRF, тонкий сервис}{%
\begin{enumerate}
  \item EDA: топы, цена vs позиция, heatmaps --- картинки в figures/.
  \item Словари брендов и категорий $\to$ artifacts/*.txt.
  \item Правила на ATTR + авто-BIO.
  \item CRF: максимизируем $P(y\mid x)$; $y$ --- теги, $x$ --- слова запроса.
  \item /extract (+ debug), чтобы любой вбил запрос и увидел JSON.
\end{enumerate}
}{%
\begin{enumerate}
  \item EDA закрыт: ноутбуки + графики.
  \item MVP жив: запрос $\to$ факты, latency в SLA.
  \item Честно: метрики пока против weak labels, не против gold.
\end{enumerate}
}

\begin{center}
\begin{tcolorbox}[colback=sand!45, colframe=mvred, boxrule=0.75pt, width=0.92\textwidth, arc=2pt, before skip=2pt, after skip=10pt]
\centering\textbf{Gate дня 2.} Нет JSON к вечеру --- эмбеддинги подождут.\\[2pt]
Сначала факты, потом красивости.
\end{tcolorbox}
\end{center}

\daysection{День 03 --- Своя разметка и шесть ударов по качеству}
Перестаём слепо верить словарю и начинаем мерить правду.

\dayblock{03}{Gold / синтетика / переобучение}{%
\begin{enumerate}
  \item Руками разметить пачку запросов (лучше сотни). Span + тип.
  \item Чуть синтетики: <<категория бренд attr>> + шум.
  \item Прогнать CRF по gold --- честный F1.
  \item Закрыть шесть пунктов ниже хотя бы в v1.
  \item Переобучить на смеси weak + gold.
\end{enumerate}
}{%
\begin{enumerate}
  \item Gold-файл в репо + скрипт оценки.
  \item Табличка <<было / стало>> по косякам.
\end{enumerate}
}

\subsection*{Шесть пунктов --- backlog, не мечты в стол}
\begin{enumerate}
  \item Факты точнее --- меньше мусора в фильтрах.
  \item Бренд и модель не путаем.
  \item Нормализация: айфон / iphone / Apple $\to$ один канон.
  \item Устойчивость к шуму: опечатки, регистр, 16гб.
  \item Своя разметка + синтетика --- честные метрики.
  \item Score пары query$\leftrightarrow$click без сказки про 100\% идеал.
\end{enumerate}

\daysection{День 04 --- Standalone: NER внизу, поиск сверху}
Эмбеддинги не вместо NER. Вместе.

\dayblock{04}{Приложение + поиск по каталогу}{%
\begin{enumerate}
  \item UI: поиск, подсветка сущностей, JSON, latency vs 100 мс.
  \item Индекс по sku\_desc. На старте ок TF-IDF+SVD.
  \item Гибрид: буст по фактам NER + близость запрос--SKU.
  \item 5--10 выдач <<до / после>> --- лучший слайд защиты.
  \item Заготовка high-score пар query--click.
\end{enumerate}
}{%
\begin{enumerate}
  \item Демо: запрос $\to$ факты + топ SKU.
  \item Схема: query $\to$ NER $\to$ facts $\to$ retrieve $\to$ SKU.
\end{enumerate}
}

\begin{center}
\begin{tikzpicture}[
  n/.style={draw=slate, fill=white, font=\small\bfseries, inner sep=5pt, rounded corners=1.5pt},
  e/.style={-{Stealth}, thick, mvred}
]
  \node[n] (q) {query};
  \node[n, right=0.5 of q] (ner) {NER / факты};
  \node[n, right=0.5 of ner] (emb) {эмбеддинги};
  \node[n, right=0.5 of emb] (sku) {top SKU};
  \draw[e] (q)--(ner); \draw[e] (ner)--(emb); \draw[e] (emb)--(sku);
\end{tikzpicture}
\end{center}

\daysection{День 05 --- Упаковка и защита}
\dayblock{05}{История для людей, не только git log}{%
\begin{enumerate}
  \item Питч: проблема $\to$ факты $\to$ метрики $\to$ демо $\to$ дальше.
  \item Сценарии: категория, бренд+ATTR, грязный запрос.
  \item Ретро без понтов: что срезали и где клики врали.
\end{enumerate}
}{%
\begin{enumerate}
  \item Защита с живым демо и честной зоной роста.
  \item After: dual-encoder, больше gold, когда-нибудь A/B.
\end{enumerate}
}

\section{Кто за что}
\secrule
\begin{tabular}{@{}p{3.1cm}p{11.6cm}@{}}
\toprule
\textbf{Зона} & \textbf{Фокус} \\
\midrule
Данные / EDA & parquet, графики, словари, score query--click \\
NER / метрики & BIO, CRF, gold, F1, нормализация \\
Поиск / app & UI, индекс, эмбеддинги, гибридная выдача \\
Доки / питч & LaTeX, README, слайды \\
\bottomrule
\end{tabular}

\vspace{0.35cm}
\section{Чего сознательно не делаем}
\secrule
\begin{enumerate}
  \item Не тащим огромный LLM в runtime, если он убивает SLA.
  \item Не рисуем 40 графиков вместо gold на день 3.
  \item Не обещаем идеальный поиск М.Видео за неделю --- обещаем контур факты$\to$выдача.
\end{enumerate}

\vspace{0.45cm}
\begin{center}
\color{inksoft}\small
Если день 2 горит --- режем эмбеддинги до TF-IDF.\\
JSON-MVP и честность про метрики не режем.
\end{center}

\end{document}
